% 第 4 章 Multipliers with Singularities (原书第 75-88 页)
\chapter{带奇性的乘子}
\label{chap:4}

在本章中, 我们将推广第~\ref{chap:2} 章的标准模型. 该模型告诉我们: 若 $\mu=\sum\mu_j\phi_j$, 其中 $\sum|\mu_j|<\infty$, 则 $\mu$ 是所有 $L^p$ 的乘子. 在 Hilbert 变换的情形, $\sum|\mu_j|=\infty$; 然而 Hilbert 变换却是 $1<p<\infty$ 时 $L^p$ 的乘子. 本章的目的是把我们对 Hilbert 变换的理解推广到更一般的乘子.

\section{H\"ormander--Mikhlin 定理}
\label{sec:4.1}
\index{H\"ormander--Mikhlin 定理 (H\"ormander--Mikhlin theorem)|emph}

\begin{proposition}\label{prop:4.1}
设 $\mu\in L^\infty$ 且 $\mu'\in L^1$. 则 $\mu$ 是 $L^p(\mathbb{R}^1)$ ($1<p<\infty$) 的乘子.
\end{proposition}

\begin{proof}
设 $\psi$ 为光滑函数, 在 $[-N,N]$ 上 $\psi=1$, 在 $[-(N+1),(N+1)]$ 之外 $\psi=0$. 令 $\mu_0=\mu\psi$; 我们将证明 ${}_p\|T_0\|\leq C$, 其中 $C$ 与 $N$ 无关. 由一个标准的极限论证, 这蕴含 $\mu$ 是乘子. 照例, 我们使用对偶技巧: 设 $f,g\in\mathcal{S}$, $\|f\|_p=\|g\|_{p'}=1$. 则
\begin{align*}
\int T_{0}fg&=\int\mu_{0}(\xi)\hat{f}(\xi)\hat{g}(\xi)\,d\xi
=\int_{-\infty}^{\infty}\int_{-\infty}^{\xi}\mu_{0}'(\eta)\,d\eta\,\hat{f}(\xi)\hat{g}(\xi)\,d\xi \\
&=\int_{-\infty}^{\infty}\int_{-\infty}^{\infty}\mu_{0}'(\eta)\,\chi_{(-\infty,\xi)}(\eta)\,\hat{f}(\xi)\hat{g}(\xi)\,d\xi\,d\eta \\
&=\int_{-\infty}^{\infty}\int_{-\infty}^{\infty}\chi_{(\eta,\infty)}(\xi)\,\hat{f}(\xi)\hat{g}(\xi)\,d\xi\,\mu_{0}'(\eta)\,d\eta.
\end{align*}
分部积分中不出现边界项, 因为 $\mu_0$ 在 $\infty$ 处消失; 由于 $\mu_0,\mu_0'$ 属于 $L^1$ 且 $\hat{f},\hat{g}$ 属于 $\mathcal{S}$, 交换积分的技术细节是容易的. 现在设 $T_\eta$ 是以
\[
\chi_{(\eta,\infty)}(\xi)=\frac{1}{2}\bigl(1-i\,\mathrm{sign}(\xi-\eta)\bigr)
\]
为乘子的算子, 于是 $T_{\eta}=\frac{1}{2}\bigl(I-iH_{\eta}\bigr)$, 其中 $H_{\eta}$ 是 Hilbert 变换的一个平移. 由此诸 $T_{\eta}$ 在 $L^{p}$ 上一致有界, 于是我们可以这样完成估计:
\begin{align*}
\int T_{0}fg&=\int_{-\infty}^{\infty}T_{\eta}f(x)\,g(x)\,dx\,\mu_{0}'(\eta)\,d\eta \\
&\leq\int\|T_{\eta}\|_{p}\,\|g\|_{p'}\,|\mu_{0}'(\eta)|\,d\eta
\leq C\|f\|_{p}\,\|g\|_{p'}\int|\mu_{0}'(\eta)|\,d\eta.
\end{align*}
为完成证明, 注意
\[
\|\mu_{0}'\|_{1}\leq\|\mu'\psi\|_{1}+\|\mu\psi'\|_{1}\leq C\|\mu'\|_{1}+\|\mu\|_{\infty}\|\psi'\|_{1}.
\]
我们只须选取 $\psi$ 使 $\|\psi'\|_{1}\leq C$, 其中 $C$ 与 $N$ 无关. 这是可能的, 因为 $\psi'$ 仅在 $(-(N+1),-N)$ 与 $(N,N+1)$ 上非零.
\end{proof}

\begin{remark}\label{rem:4.2}
我们稍稍改变了视角; 我们不再考察 $\sum\mu_j\psi_j$ (其中诸 $\psi_j$ 紧支集), 而是考察 Hilbert 变换的加权平均. 这两种看法是一样的, 正如积分是和的连续版本.

牢记这一点, 条件 $\mu'\in L^1$ 便类似于条件 $\sum|\mu_j|<\infty$. 简言之, 除了记号之外, 该命题并没有真正使我们跳出标准模型.

我们真正想知道的, 是条件``$\mu=\sum\psi_j$ 而前面没有收敛因子''的连续类比. 这不难看出; 直观上, $\psi_j'=2^j$ 于区间 $[2^{-(j+1)},2^{-j}]$ 上. 用 $\mu$ 的语言, $\mu'(\xi)=\frac{1}{\xi}$. 正确的乘子条件是 $\xi\mu'(\xi)\in L^{\infty}$.
\end{remark}

\begin{theorem}\label{thm:4.3}
设 $\mu\in L^{\infty}$ 且 $\mu$ 在 $0$ 附近之外是 $C^{1}$ 的. 若 $\xi\mu'(\xi)\in L^{\infty}$, 则 $\mu$ 是 $L^{p}$ ($1<p<\infty$) 的乘子.
\end{theorem}

\begin{proof}
我们打算应用 Hilbert 变换理论; 我们需要把二进分解 $\mu=\sum\psi_j$ 转移到对卷积核的估计上. 为此, 对 $\mu$ 的分解必须光滑地进行, 使得逆 Fourier 变换把 $\mathcal{S}$ 函数映到 $\mathcal{S}$ 函数. 技术细节: 设 $\varphi\geq 0$ 光滑, 且满足
\[
\varphi(\xi)=1\quad\text{若}\ \frac{3}{4}<\xi<\frac{3}{2},
\]
\[
\varphi(\xi)=0\quad\text{若}\ \xi<\frac{1}{2}\ \text{或}\ \xi>2.
\]
令
\[
\psi_{j}(\xi)=\frac{\varphi\bigl(2^{-j}|\xi|\bigr)}{\sum_{k=-\infty}^{\infty}\varphi\bigl(2^{-k}|\xi|\bigr)}.
\]
再令
\[
\psi_{j}(\xi)=\psi_{0}\bigl(2^{-j}\xi\bigr);
\]
设 $\mu_j=\psi_j\mu$, 并定义乘子三元组 $(T,K,\mu)$ 与 $(T_j,K_j,\mu_j)$. 我们的希望是应用 Hilbert 变换的思想, 于是我们希望
\[
\int_{|x|>2|y|}|K_{j}(x-y)-K_{j}(x)|\,dx\leq C2^{-j}.
\]
不确定性原理阻止了这一点; 由于 $\mu_{j}$ 支在长度为 $2^{j}$ 的区间上, $K_{j}$ 直观上集中在一个长度为 $2^{-j}$ 的区间上. 因此, 若 $|y|>2^{-j}$, 则 $K_{j}(x-y)$ 绝不会看起来像 $K_{j}(x)$, 从而我们无法利用光滑性带来的控制. 我们唯一的希望是 $K_{j}$ 满足估计:
\begin{align*}
\int_{|x|>1}|K_{j}(x)|\,dx
&\leq C\left(\int_{|x|>1}|x|^{2}|K_{j}(x)|\,dx\right)^{\frac{1}{2}}
\leq C\|xK_{j}\|_{2}=C\Bigl\|\frac{d}{d\xi}\hat{K}_{j}\Bigr\|_{2}=C\|\mu_{j}'\|_{2}.
\end{align*}
但是
\begin{align*}
\|\left(\psi_{j}\mu\right)^{\prime}\|_{2}&\leq\|\psi_{j}\mu^{\prime}\|_{2}+\|\psi_{j}^{\prime}\mu\|_{2}, \\
\int\bigl|D^{1-k}\psi_{j}\,D^{k}\mu\bigr|^{2}
&=\int_{\operatorname{supp}\psi_{j}}\left|2^{-j(1-k)}\psi_{0}^{(1-k)}\bigl(2^{-j}|\xi|\bigr)\,D^{k}\mu\right|^{2} \\
&\leq C\,2^{-2(1-k)j}\int_{\operatorname{supp}\psi_{j}}\bigl|D^{k}\mu\bigr|^{2}
\leq C\,2^{-2(1-k)j}\int_{2^{j-1}}^{2^{j+1}}|\xi|^{-2k}\,d\xi
=C\bigl(2^{-\frac{j}{2}}\bigr)^{2},
\end{align*}
因此 $\|xK_j\|_2\leq C2^{-\frac{j}{2}}$. 当无法使用光滑性时, 我们将应用这类估计; 于是对 $2^j|y|\geq 1$,
\[
\int_{|x|>2|y|}|K_{j}(x-y)-K_{j}(x)|\,dx
\leq\int_{|x|>2|y|}|K_{j}(x-y)|\,dx+\int_{|x|>2|y|}|K_{j}(x)|\,dx.
\]
在第一个积分中, 变量替换把我们带到集合
\[
\{x \mid |x+y|>2|y|\}\subset\{|x|+|y|>2|y|\}=\{|x|>|y|\},
\]
而两个积分都被
\begin{align*}
\int_{|x|>|y|}|K_{j}(x)|\,dx
&\leq\left(\int_{|x|>|y|}|x|^{-2}\,dx\right)^{\frac{1}{2}}\|xK_{j}\|_{2} \\
&\leq C|y|^{-\frac{1}{2}}\,2^{-\frac{j}{2}}=C\bigl(2^{j}|y|\bigr)^{-\frac{1}{2}}
\end{align*}
控制. 接着, 在区域 $|y|<2^{-j}$ 中, 我们预计可以使用光滑性. 若 $|y|<2^{-j}$ 且 $|x|>2|y|$, 则
\begin{align*}
\int_{\{|x|>2|y|\}}|K_{j}(x-y)-K_{j}(x)|\,dx
&\leq\left(\int_{|x|\leq 2^{-j}}1\,dx\right)^{\frac{1}{2}}\|K_{j}(x-y)-K_{j}(x)\|_{2} \\
&=2^{-\frac{j}{2}}\bigl\|\bigl(e^{2\pi iy\xi}-1\bigr)\psi_{j}(\xi)\mu(\xi)\bigr\|_{2}
\leq C2^{-\frac{j}{2}}\|\xi\psi_{j}\mu\|_{2}\,|y|.
\end{align*}
在 $\psi_j$ 的支集上, $|\xi|\leq 2^j$ 且 $|\psi_j\mu|\leq C$, 所以这一项被
\[
C2^{-\frac{j}{2}}\,2^{j}|y|\,|\operatorname{supp}\psi_{j}|^{\frac{1}{2}}=C\bigl(2^{j}|y|\bigr)
\]
控制. 最后, 若 $|x|>2^{-j}$,
\begin{align*}
\int_{|x|>2^{-j}}\left|K_{j}(x-y)-K_{j}(x)\right|dx
&\leq\left(\int_{|x|>2^{-j}}|x|^{-2}\,dx\right)^{\frac{1}{2}}\bigl\|x\bigl[K_{j}(x-y)-K_{j}(x)\bigr]\bigr\|_{2} \\
&=\;C2^{\frac{j}{2}}\bigl\|\bigl(e^{2\pi iy\xi}-1\bigr)\bigl(\psi_{j}(\xi)\mu(\xi)\bigr)^{\prime}\bigr\|_{2} \\
&\qquad+2^{-\frac{j}{2}}\bigl\|\bigl(e^{2\pi iy\xi}-1\bigr)\psi_{j}(\xi)\mu(\xi)\bigr\|_{2} \\
&\leq C2^{\frac{j}{2}}|y|\,\bigl\|\xi(\psi_{j}\mu)^{\prime}\bigr\|_{2}+C2^{\frac{j}{2}}|y|\,\|\psi_{j}\mu\|_{2}.
\end{align*}
在第一项中, 在 $\psi_j$ 的支集上 $|\xi|\leq 2^j$, 于是
\[
\bigl\|\bigl(\psi_{j}\mu\bigr)^{\prime}\bigr\|_{2}\leq 2^{-\frac{j}{2}}.
\]
在第二项中,
\[
\|\psi_{j}\mu\|_{2}\leq\|\mu\|_{\infty}\|\psi_{j}\|_{2}=C2^{\frac{j}{2}},
\]
而且两项都被 $C2^{j}|y|$ 控制.

为完成证明, 设 $\mu_N=\sum_{j=-N}^{N}\psi_j\mu$, 并设 $(T_N,K_N,\mu_N)$ 为相应的乘子三元组. 由于
\[
\sum_{j=-N}^{N}\psi_{j}\leq 1,\qquad\|\mu_{N}\|_{\infty}\leq\|\mu\|_{\infty},
\]
而且我们有主要估计
\[
\int_{|x|>2|y|}|K_{N}(x-y)-K_{N}(x)|\,dx
\leq C\sum_{2^{j}|y|\leq 1}2^{j}|y|+C\sum_{2^{j}|y|>1}\bigl(2^{j}|y|\bigr)^{-\frac{1}{2}}.
\]
这两个和关于 $y$ 一致有界; 例如若 $|y|>1$, 则第一个和中的指数 $j$ 全是负的, 其贡献为
\[
C|y|\sum_{j\leq-\log|y|}2^{j}=C|y|\,|y|^{-1},
\]
而在第二项中,
\[
|y|^{-\frac{1}{2}}\sum_{j\geq-\log|y|}2^{-\frac{j}{2}}
\leq C\bigl(1+|y|^{-\frac{1}{2}}\bigr).
\]
但 $|y|^{-\frac{1}{2}}\leq 1$. 对 $|y|\leq 1$ 有类似的估计.

现在只需对 $(T_N,K_N,\mu_N)$ 应用定理~\ref{thm:2.32}, 得到它们在 $L^p$ 上有界, 且界与 $N$ 无关. 一个快速的极限论证告诉我们如何把算子 $T$ 扩张到整个 $L^p$: 若 $f\in\mathcal{S}$, 则 $\sum\psi_j\leq 1$ 且 $\mu\in L^\infty$, 于是 $\mu_N\hat{f}$ 控制收敛且在 $L^2$ 中收敛到 $\mu\hat{f}$. 算子 $T_N$ 一致有界且在一个稠密集上收敛到 $T$, 故 $T$ 也有界. 在 $\mathcal{S}$ 上, $T$ 由 $\mu$ 给出.
\end{proof}

\begin{remark}\label{rem:4.4}
\begin{enumerate}[label=\alph*)]
\item Hilbert 变换及其在命题~\ref{prop:4.1} 中的简单变体, 证明了具有简单跳跃间断的乘子的有界性. 我们新的、改进了的结果涵盖了一些振荡型的间断; 例如 $\mu(\xi)=|\xi|^{it}$ 满足定理的假设.
\item 有时, 检验较弱的条件更容易:
\[
\sup_{R>0}\frac{1}{R^{1+2j}}\int_{R<|\xi|<2R}\bigl|D^{j}\mu\bigr|^{2}\leq C,\qquad j=0,\ 1.
\]
在定理~\ref{thm:4.3} 的证明中, 这实际上就是我们真正需要的全部.
\end{enumerate}
\end{remark}

\essayhead{ess:4.5}{关于标准模型的注记}

我们刚完成的证明只在 $f_j$ 的支集上估计 $D^j\mu$, 而我们取后者为一个二进区间. 在这个意义上, 我们把 $\mu_j=\mu\psi_j$ 视为彼此独立地行动的乘子. 在本章的其余部分, 我们将把这一想法严格化. 如同上面的证明, 我们将给出三个版本. 首先考察没有光滑性的乘子, 比如 Hilbert 变换:
\[
\widehat{\left(P_{j}f\right)}(\xi)=\chi_{(2^{j},2^{j+1})}(\xi)\,\hat{f}(\xi).
\]
接下来我们证明光滑版本, 其中 $\varphi\in C^\infty$ 且 $\operatorname{supp}\varphi\subset[1,2]$; 此时
\[
\widehat{\left(T_{j}f\right)}(\xi)=\varphi(2^{-j}\xi)\,\hat{f}(\xi).
\]
最后是乘子版本:
\[
\widehat{\left(T_{j}f\right)}(\xi)=\mu(\xi)\,\varphi(2^{-j}\xi)\,\hat{f}(\xi),
\]
其中 $\mu$ 满足与~\ref*{thm:4.3} 中类似的估计. 与之前一样, 我们将用不连续的 $P_j$ 作为证明乘子定理的技术工具; 这最终将消除对 $K$ 的核估计, 使我们能完全集中在 $\mu$ 上. 最终的好处是我们可以完全用标准模型中的想法工作, 而不必在 $T$ 与 $\mu$ 的估计之间来回切换.

这个纲领基于对``诸 $\mu_j$ 算子彼此独立地行动''的含义的理解. 若转而考察以 $\chi_{(2^j,2^{j+1})}$ 为乘子的算子 $P_j$, 独立性表现为这些算子是正交投影这一事实:
\[
\|f\|_{2}^{2}=\|\hat{f}\|_{2}^{2}=\sum_{j=-\infty}^{\infty}\int_{2^{j}}^{2^{j+1}}|\hat{f}|^{2}
=\sum\int|P_{j}f|^{2}=\int\Bigl(\sum|P_{j}f|^{2}\Bigr).
\]
用 $L^{2}$ 估计的语言, 独立性的正确表达是:
\[
\|f\|_{2}=\Bigl\|\Bigl(\sum|P_{j}f|^{2}\Bigr)^{\frac{1}{2}}\Bigr\|_{2}.
\]
现在我们必须猜测``独立性''或``正交性''在 $L^p$ 上的类比; 我们的出发点是量 $\bigl(\sum|P_jf|^2\bigr)^{\frac{1}{2}}$, 只要并非所有的 $|P_jf|$ 同时很大, 它就像 $|P_jf|$ 一样. 定理将给出: 这就是正交性的正确观念:
\[
\Bigl\|\Bigl(\sum|P_{j}f|^{2}\Bigr)^{\frac{1}{2}}\Bigr\|_{p}\leq C_{p}\|f\|_{p}.
\]

\begin{lemma}\label{lem:4.6}
设
\[
\Bigl\|\Bigl(\sum|P_{j}f|^{2}\Bigr)^{\frac{1}{2}}\Bigr\|_{p}\leq C_{p}\|f\|_{p}
\]
对所有的 $f\in L^{p}$ 成立. 则对所有的 $g\in L^{p'}$,
\[
\|g\|_{p'}\leq C_{p}\Bigl\|\Bigl(\sum|P_{j}g|^{2}\Bigr)^{\frac{1}{2}}\Bigr\|_{p'}.
\]
\end{lemma}

\begin{proof}
\begin{align*}
\int fg&=\int\hat{f}\,\check{g}=\sum\int\widehat{P_{j}f}\,\widehat{P_{j}g} \\
&=\int\sum\bigl(P_{j} f\bigr)\bigl(P_{j} g\bigr) \\
&\leq\int\Bigl(\sum|P_{j} f|^{2}\Bigr)^{\frac{1}{2}}\Bigl(\sum|P_{j} g|^{2}\Bigr)^{\frac{1}{2}} \\
&\leq\Bigl\|\Bigl(\sum|P_{j} f|^{2}\Bigr)^{\frac{1}{2}}\Bigr\|_{p}\Bigl\|\Bigl(\sum|P_{j} g|^{2}\Bigr)^{\frac{1}{2}}\Bigr\|_{p'} \\
&\leq C_{p}\|f\|_{p}\Bigl\|\Bigl(\sum|P_{j} g|^{2}\Bigr)^{\frac{1}{2}}\Bigr\|_{p'}.
\end{align*}
现在对所有满足 $\|f\|_p=1$ 的 $f$ 取上确界, 即得结论.
\end{proof}

\begin{remark}\label{rem:4.7}
\index{独立性 (independence)}
这个结果的全部要点在于下述推论: 当我们对 $1<p<\infty$ 证明了
\[
\Bigl\|\Bigl(\sum|P_{j}f|^{2}\Bigr)^{\frac{1}{2}}\Bigr\|_{p}\leq C_{p}\|f\|_{p}
\]
时, 我们自动得到
\[
\|f\|_{p}\leq C_{p}\Bigl\|\Bigl(\sum|P_{j}f|^{2}\Bigr)^{\frac{1}{2}}\Bigr\|_{p}\leq C_{p}\|f\|_{p}.
\]
为了比较 $f$ 与 $\bigl(\sum|P_jf|^2\bigr)^{1/2}$ 的 $L^p$ 范数, 我们需要几个新想法; 其复杂性来自这样一个事实: 我们是在为 $p\neq 2$ 的 Banach 空间 $L^p$ 定义一种``正交性''的观念, 而在那里正交性根本没有任何意义. 正确的想法来自概率论, 称为独立性 (independence). 若 $f$, $g$ 是可测函数, 我们说它们独立, 如果对一切可测集 $E$ 有 $\int_Efg=\int_Ef\int_Eg$; 其想法是: $f$ 的大小必须在逐点意义上影响 $g$ 的大小, 但在小集合上取平均时, 两个函数互不影响. 若用二项式定理把一个和展开到非整数次幂, 你就必须估计的正是这类对象.
\end{remark}

\section{Littlewood--Paley 与 Marcinkiewicz 定理}
\label{sec:4.2}

\begin{definition}\label{def:4.8}
\index{Rademacher 函数 (Rademacher functions)}
Rademacher 函数 $r_{m}(t)$ 由以下要求定义: 它们以 $2$ 为周期, 且
\begin{gather*}
r_{0}(t)=r_{0}(t+1); \\
r_{0}(t)=1\quad\text{若}\ 0\leq t\leq\frac{1}{2}; \\
r_{0}(t)=-1\quad\text{若}\ \frac{1}{2}<t\leq 1; \\
r_{m}(t)=r_{0}(2^{m}t).
\end{gather*}
\end{definition}

\begin{theorem}\label{thm:4.9}
Rademacher 函数在 $[0,1]$ 上独立; 也就是说, 对 $m_1,m_2,\ldots,m_k$ 与 $\alpha_1,\allowbreak\alpha_2,\ldots,\alpha_k$,
\[
|\{t \mid r_{m_{i}}(t)\leq\alpha_{i}\ \text{对所有的}\ i\}|=\prod_{i=1}^{k}|\{t \mid r_{m_{i}}(t)\leq\alpha_{i}\}|,
\]
而且对可测函数 $f_i$ 与任意可测集 $E$,
\[
\int_{E}\prod_{i}f_{i}\bigl(r_{m_{i}}(t)\bigr)\,dt=\prod_{i}\int_{E}f_{i}\bigl(r_{m_{i}}(t)\bigr)\,dt.
\]
\end{theorem}

\begin{proof}
关于积分的结果源于关于阶梯函数的类似结果, 后者源于关于 Borel 集的结果, 再源于关于区间的结果, 而后者又源于关于半无穷区间的结果, 我们就来证明它. 我们只处理 $(-\infty,\alpha_i]$ 的情形.

若对某个 $i$ 有 $\alpha_i<-1$, 则两边都有一些集合是空的, 乘积为零. 若对某个 $i$ 有 $\alpha_i>1$, 则条件 $r_{m_i}(t)\leq\alpha_i$ 总被满足, 于是它在左端没有贡献, 而在右端贡献因子 $1$. 类似地, 条件 $r_m(t)\leq\alpha_m<1$ 等价于 $r_m(t)\leq-1$, 即 $r_m(t)=-1$. 于是, 定理由等式
\[
|\{t \mid r_{m_{i}}(t)=-1\ \text{对所有的}\ i\}|
=\prod_{i=1}^{k}|\{t \mid r_{m_{i}}(t)=-1\}|=\left(\frac{1}{2}\right)^{k}
\]
得出.

这样的计算是一个归纳法, 其中除最后一步外都被跳过了. 我们重新排序使 $m_1<m_2<\cdots<m_k$. 归纳假设是: $\{t \mid r_{m_i}(t)=-1,\ 1\leq i\leq k-1\}$ 是二进区间的无交并, 每个的长度为 $2^{-m_{k-1}}$, 而该集合的总测度为 $\bigl(\frac{1}{2}\bigr)^{k-1}$. 现在再加上条件 $r_{m_k}(t)=-1$. 由于 $m_k>m_{k-1}$, $r_{m_k}$ 在这些二进区间的每一个上恰在一半的时间里取负值 (为什么? 你必须用到 $r_{m_k}$ 是 $r_{m_{k-1}}$ 的伸缩这一事实).
\end{proof}

\begin{theorem}\label{thm:4.10}
设 $F(t)=\sum_{k=0}^{N}a_kr_k(t)$. 若 $\sum a_k^2=1$, 则
\[
|\{t \mid |F(t)|>\lambda\}|\leq 2e^{-\frac{\lambda^{2}}{4}}.
\]
\end{theorem}

\begin{proof}
\begin{align*}
e^{\frac{\lambda^{2}}{2}}\left|\left\{t \mid |F(t)|>\lambda\right\}\right|
&\leq\int_{|F|>\lambda}e^{\frac{\lambda}{2}|F(t)|}\,dt
\leq\int_{0}^{1}e^{\frac{\lambda}{2}|F|}\,dt \\
&\leq\int_{0}^{1}\left(e^{\frac{\lambda}{2}F}+e^{-\frac{\lambda}{2}F}\right)dt.
\end{align*}
但是
\begin{align*}
\int_{0}^{1}e^{-\frac{\lambda}{2}F}\,dt
&=\int_{0}^{1}\prod_{k}e^{-\frac{\lambda}{2}a_{k}r_{k}}\,dt
=\prod_{k}\int_{0}^{1}e^{-\frac{\lambda}{2}a_{k}r_{k}}\,dt
=\prod_{k}2^{-k}\int_{0}^{2^{k}}e^{-\frac{\lambda}{2}a_{k}r_{0}}\,dt \\
&=\prod_{k=1}^{N}\left[\sum_{n=0}^{2^{k-1}}\left(\int_{n}^{n+\frac{1}{2}}e^{-\frac{\lambda}{2}a_{k}}\,dt+\int_{n+\frac{1}{2}}^{n+1}e^{\frac{\lambda}{2}a_{k}}\,dt\right)\right] \\
&=\prod_{k=1}^{N}2^{-k}\left(\frac{e^{\frac{\lambda}{2}a_{k}}+e^{-\frac{\lambda}{2}a_{k}}}{2}\right)\sum_{n=0}^{2^{k-1}}1 \\
&=\prod_{k=1}^{N}\cosh\Bigl(\frac{\lambda}{2}a_{k}\Bigr)
\leq\prod_{k=1}^{N}e^{\frac{\lambda^{2}}{4}a_{k}^{2}}
=e^{\frac{\lambda^{2}}{4}\sum a_{k}^{2}}=e^{\frac{\lambda^{2}}{4}}.
\end{align*}
对另一项有类似的估计, 这带来了因子 $2$.
\end{proof}

\begin{corollary}\label{cor:4.11}
设 $F(t)=\sum a_kr_k(t)$. 若 $1<p<\infty$, 则
\[
\|F\|_{p}\leq C_{p}\|F\|_{2}\leq C_{p}'\|F\|_{p}.
\]
\end{corollary}

\begin{proof}
令
\[
G(t)=\|F\|_{2}^{-1}F(t)=\left(\sum a_{k}^{2}\right)^{-\frac{1}{2}}F(t).
\]
则 $G=\sum b_kr_k$, 其中
\[
\|G\|_{2}=\Bigl(\sum b_{k}^{2}\Bigr)^{\frac{1}{2}}=1.
\]
因此,
\begin{align*}
\|G\|_{p}^{p}&=p\int_{0}^{\infty}\lambda^{p-1}|\{|G|>\lambda\}|\,d\lambda
\leq 2p\int_{0}^{\infty}\lambda^{p-1}e^{-\frac{\lambda^{2}}{4}}\,d\lambda=C_{p},
\end{align*}
于是对 $1\leq p<\infty$, $\|F\|_p\leq C_p\|F\|_2$. 重复引理~\ref{lem:4.6} 的证明, 我们得到 $\|F\|_2^2\leq\|F\|_p\|F\|_r$, 其中 $\frac{1}{2}=\frac{1}{2r}+\frac{1}{2p}$. 若 $p>1$, 则 $r<\infty$, 于是 $\|F\|_r\leq C_p\|F\|_2$.
\end{proof}

\noindent\textbf{注:} \index{向量值不等式 (vector-valued inequalities)}这就是独立性的要点所在: $L^p$ 范数与 $L^2$ 范数等价. $L^2$ 是 Hilbert 空间, 在那里我们可以使用正交性; 该推论使我们能把它转移到 $L^p$. 原始的证明是用二项式定理把和展开成幂, 再利用独立性把乘积的积分化成积分的乘积. 通常这只能在 $L^2$ 上借助那里的正交性完成; 独立性提供了一个替代品.

\begin{proposition}\label{prop:4.12}
设 $\{I_{j}\}$ 是区间的任意集合族, 并设
\[
\widehat{\left(S_{j}f\right)}(\xi)=\chi_{I_{j}}(\xi)\,\hat{f}(\xi).
\]
则对 $1<p<\infty$,
\[
\Bigl\|\Bigl(\sum|S_{j}f_{j}|^{2}\Bigr)^{\frac{1}{2}}\Bigr\|_{p}\leq C_{p}\Bigl\|\Bigl(\sum|f_{j}|^{2}\Bigr)^{\frac{1}{2}}\Bigr\|_{p}.
\]
\end{proposition}

\begin{proof}
由于区间的特征函数可以由 Hilbert 变换构造出来, 我们先对它们证明一个类似的结果. 若 $p>2$, 令 $r=\bigl(\frac{p}{2}\bigr)'$, 并任取 $g\in L^r$. 则对几乎所有的 $x$,
\begin{align*}
|g(x)|&=\lim_{|I|\to0}\left(\frac{1}{|I|}\int_{I}|g(y)|^{s}\,dy\right)^{\frac{1}{s}}
\leq\sup_{I}\left(\frac{1}{|I|}\int_{I}|g(y)|^{s}\,dy\right)^{\frac{1}{s}}=M_{s}(g).
\end{align*}
但定理~\ref{cor:3.16} 告诉我们 $M_{s}(g)$ 属于 $A_{1}$, 而引理~\ref{lem:3.15} 告诉我们 $M_{s}(g)$ 属于 $A_{2}$, 于是定理~\ref{thm:3.13} 与命题~\ref{prop:3.14} 表明
\[
\int|H f|^{2}\,M_{s}(g)\leq C\int|f|^{2}\,M_{s}(g).
\]
因此,
\begin{align*}
\int\Bigl(\sum|H f_{j}|^{2}\Bigr)g
&\leq\int\Bigl(\sum|H f_{j}|^{2}\Bigr)M_{s}(g)
\leq C\int\Bigl(\sum|f_{j}|^{2}\Bigr)M_{s}(g) \\
&\leq C\Bigl\|\sum|f_{j}|^{2}\Bigr\|_{\frac{p}{2}}\bigl\|M_{s}(g)\bigr\|_{(\frac{p}{2})'}
=C\Bigl\|\Bigl(\sum|f_{j}|^{2}\Bigr)^{\frac{1}{2}}\Bigr\|_{p}^{2}\,\|M_{s}(g)\|_{r}.
\end{align*}
但 $r>1$, 于是我们可以找到 $s>1$ 使 $\frac{r}{s}>1$, 而推论~\ref{cor:2.25} 给出估计
\begin{align*}
\|M_{s}(g)\|_{r}^{r}&=\int\bigl[(|g|^{s})^{\star}\bigr]^{\frac{r}{s}}
\leq C\int\bigl[|g|^{s}\bigr]^{\frac{r}{s}}=C\int|g|^{r}=C\|g\|_{r}^{r}.
\end{align*}
总之,
\[
\int\Bigl(\sum|H f_{j}|^{2}\Bigr)g\leq C\Bigl\|\Bigl(\sum|f_{j}|^{2}\Bigr)^{\frac{1}{2}}\Bigr\|_{p}^{2}\|g\|_{r}.
\]
若对所有满足 $\|g\|_r=1$ 的 $g$ 取上确界, 我们便在左端得到 $r'=\frac{p}{2}$ 范数, 而这正是
\[
\left(\int\Bigl(\sum|H f_{j}|^{2}\Bigr)^{\frac{p}{2}}\right)^{\frac{2}{p}}
=\Bigl\|\Bigl(\sum|H f_{j}|^{2}\Bigr)^{\frac{1}{2}}\Bigr\|_{p}^{2}.
\]
在 $p=2$ 的情形,
\begin{align*}
\Bigl\|\Bigl(\sum|H f_{j}|^{2}\Bigr)^{\frac{1}{2}}\Bigr\|_{2}^{2}
&=\int\sum|H f_{j}|^{2}=\int\sum|f_{j}|^{2}
=\Bigl\|\Bigl(\sum|f_{j}|^{2}\Bigr)^{\frac{1}{2}}\Bigr\|_{2}^{2},
\end{align*}
因为 $H$ 是 $L^{2}$ 上的等距算子. $1<p<2$ 的情形用对偶性完成, 正如引理~\ref{lem:4.6} 的证明那样.

接下来我们需要从 $H$ 过渡到区间. 设 $I_{j}=[a_{j},b_{j}]$; 则
\[
S_{j}=\frac{1}{4}\bigl(I+iH_{a_{j}}\bigr)\bigl(I-iH_{-b_{j}}\bigr),
\]
正如引理~\ref{lem:2.12} 中那样. 使用定理~\ref{thm:1.23}d) 的记号, $H_{c}f=J(Hf)J^{-1}$, 其中 $Jf=e^{2\pi icx}f$. 由此, $\bigl(I+iH_{a_{j}}\bigr)\bigl(I-iH_{-b_{j}}\bigr)$ 可以由四个算子控制, 其中每一个都由 $H$ 或 $I$ 控制.
\end{proof}

\begin{theorem}\label{thm:4.13}
\index{Littlewood--Paley 分解 (Littlewood--Paley decomposition)}
设 $\Delta_j=(2^j,2^{j+1})\cup(-2^{j+1},-2^j)$, $-\infty<j<\infty$. 设 $P_j$ 是以 $\chi_{\Delta_j}$ 为乘子的乘子算子. 则对 $1<p<\infty$,
\[
\|f\|_{p}\leq C_{p}\Bigl\|\Bigl(\sum|P_{j}f|^{2}\Bigr)^{\frac{1}{2}}\Bigr\|_{p}\leq C_{p}'\|f\|_{p}.
\]
\end{theorem}

\begin{proof}
在 $p=2$ 处成立的是精确等式, 于是引理~\ref{lem:4.6} 告诉我们, 只需证明右端的不等式. 我们将使用给出有界乘子的 $\chi_{(-\infty,0)}$, 分别构造 $(2^j,2^{j+1})$ 及其负部. 我们的策略是用命题~\ref{prop:4.12} 比较
\[
\Bigl\|\Bigl(\sum|P_{j}f|^{2}\Bigr)^{\frac{1}{2}}\Bigr\|_{p}
\]
与
\[
\Bigl\|\Bigl(\sum|T_{j}f|^{2}\Bigr)^{\frac{1}{2}}\Bigr\|_{p},
\]
其中诸 $T_{j}$ 是光滑的. 然后我们将用推论~\ref{cor:4.11} 比较
\[
\Bigl\|\Bigl(\sum|T_{j}f|^{2}\Bigr)^{\frac{1}{2}}\Bigr\|_{p}
\]
与 $\bigl\|\sum T_jf\bigr\|_p$. 最后, 我们将利用 Rademacher 函数把两者拼合起来. 设
\begin{gather*}
\phi(\xi)=1\quad\text{对}\ 1\leq\xi\leq 2; \\
\phi(\xi)=0\quad\text{对}\ \xi\notin\Bigl[\frac{1}{2},\frac{5}{2}\Bigr]; \\
\phi_{j}(\xi)=\phi(2^{j}\xi).
\end{gather*}
则 $\chi_{\Delta_j}=\phi_j\chi_{\Delta_j}$, 且若 $T_j$ 表示以 $\phi_j$ 为乘子的算子, 则 $P_jf=P_jT_jf$, 于是
\begin{align*}
\Bigl\|\Bigl(\sum|P_{j} f|^{2}\Bigr)^{\frac{1}{2}}\Bigr\|_{p}
&=\Bigl\|\Bigl(\sum|P_{j} T_{j} f|^{2}\Bigr)^{\frac{1}{2}}\Bigr\|_{p}
\leq C\Bigl\|\Bigl(\sum|T_{j} f|^{2}\Bigr)^{\frac{1}{2}}\Bigr\|_{p}.
\end{align*}
现在令
\[
\mu_{o}(\xi,t)=\sum_{j\ \text{奇}}\phi_{j}(\xi)\,r_{j}(t);
\]
对 $\mu_{e}$ (偶数指标) 类似. 我们将证明 $\mu_{o},\mu_{e}$ 关于 $t$ 一致地满足定理~\ref{thm:4.3} 的假设. 设 $T_{o}$, $T_{e}$ 为相应的算子. 则由推论~\ref{cor:4.11},
\[
\Bigl(\sum|T_{j}f|^{2}\Bigr)^{\frac{1}{2}}\leq C\Bigl\|\sum T_{j}f(x)\,r_{j}(t)\Bigr\|_{p,t}.
\]
于是
\begin{align*}
\int\Bigl(\sum|T_{j}f|^{2}\Bigr)^{\frac{p}{2}}dx
&\leq C\int_{0}^{1}\int\Bigl|\sum T_{j}f(x)\,r_{j}(t)\Bigr|^{p}\,dx\,dt
=C\int_{0}^{1}\int|T_{o}f(x,t)|^{p}\,dx\,dt \\
&\leq C\int_{0}^{1}\int|f(x)|^{p}\,dx\,dt=C\|f\|_{p}^{p}.
\end{align*}
这就证明了定理; 为得到关于 $\mu_{o}$ 的估计, 注意奇数的诸 $\phi_{j}$ 具有互不相交的支集, 于是
\[
|\mu_{o}(\xi,t)|\leq\sum|\phi_{j}(\xi)|\leq 1.
\]
类似地, 若 $\xi\in\operatorname{support}\phi_{j}$, 则
\begin{align*}
|\xi\,\mu_{0}'(\xi,t)|
&\leq|\xi|\,2^{-j}\bigl|\phi_{0}'\bigl(2^{-j}\xi\bigr)\bigr|
\leq 2^{j+1}\,2^{-j}\,\|\phi_{0}'\|_{\infty}\leq C. \qedhere
\end{align*}
\end{proof}

\begin{theorem}\label{thm:4.14}
\index{Marcinkiewicz 乘子定理 (Marcinkiewicz multiplier theorem)}
设 $\mu\in L^{\infty}$, 且对每个二进区间 $I$,
\[
\int_{I}|\mu^{\prime}|\leq C
\]
关于 $I$ 一致成立. 则 $\mu$ 给出 $L^{p}$ ($1<p<\infty$) 的一个乘子.
\end{theorem}

\begin{proof}
首先注意
\[
\left\|T f\right\|_{p}\leq C\Bigl\|\Bigl(\sum|P_{j}T f|^{2}\Bigr)^{\frac{1}{2}}\Bigr\|_{p}.
\]
若 $I_{j}=(2^{j},2^{j+1})$, 则
\begin{align*}
P_{j}Tf(x)&=\int_{I_{j}}\mu(\xi)\hat{f}(\xi)e^{2\pi ix\xi}\,d\xi \\
&=\int_{I_{j}}\int_{2^{j}}^{\xi}\mu'(\eta)\,d\eta\,\hat{f}(\xi)\chi_{I_{j}}(\xi)\,e^{2\pi ix\xi}\,d\xi
+\mu(2^{j})\int\hat{f}(\xi)\chi_{I_{j}}(\xi)\,e^{2\pi ix\xi}\,d\xi \\
&=\int_{I_{j}}\mu'(\eta)\int_{\eta}^{2^{j+1}}\hat{f}(\xi)\chi_{I_{j}}(\xi)\,e^{2\pi ix\xi}\,d\xi\,d\eta
+\mu(2^{j})P_{j}f(x) \\
&=\int_{I_{j}}\mu'(\eta)\,\frac{(I+iH_{\eta})}{2}\,P_{j}f(x)\,d\eta+\mu(2^{j})P_{j}f(x).
\end{align*}
由于 $|\mu(2^j)|\leq\|\mu\|_\infty$, 我们可以把第二项分离出来, 然后使用估计
\[
\Bigl\|\Bigl(\sum|P_{j}f|^{2}\Bigr)^{\frac{1}{2}}\Bigr\|_{p}\leq C\|f\|_{p}.
\]
由于
\[
\int_{I}\mu'(\eta)\,\frac{1}{2}I\bigl(P_{j}f(x)\bigr)\,d\eta
=\frac{1}{2}\left(\int_{I}\mu'\right)P_{j}f(x)\leq C\,P_{j}f(x),
\]
这一项也可以被控制. ``主要的''项总是你最后才能对付的那一项; 在此情形, $|P_jTf|^2$ 中的主要项是
\begin{align*}
\left|\int_{I_{j}}\mu'(\eta)\,H_\eta P_{j}f(x)\,d\eta\right|^{2}
&\leq\left(\int_{I_{j}}|\mu'(\eta)|\right)\int_{I_{j}}|\mu'|\,|H_\eta P_{j}f(x)|^{2}\,d\eta.
\end{align*}
如同命题~\ref{prop:4.12} 的证明, 我们取 $g\in L^r$, $r=\bigl(\frac{p}{2}\bigr)'$, 并注意到 $g\leq M_s(g)$. 则
\begin{align*}
\int\sum|P_{j}T f|^{2}\,g(x)\,dx
&\leq 2C\int\sum|P_{j}f|^{2}\,g
+C\sum\int_{I_{j}}|\mu'(\eta)|\int|H_\eta P_{j}f|^{2}\,g(x)\,dx\,d\eta \\
&\leq C\Bigl\|\Bigl(\sum|P_{j}f|^{2}\Bigr)^{\frac{1}{2}}\Bigr\|_{p}^{2}\,\|g\|_{r} \\
&\qquad+C\sum\int_{I_{j}}|\mu'(\eta)|\int|H_\eta P_{j}f|^{2}\,M_{s}(g)\,dx\,d\eta.
\end{align*}
$A_{p}$ 理论适用于诸 $H_{\eta}$ 算子; 界与 $\eta$ 无关:
\[
\int\sum|P_{j}T f|^{2}\,g(x)\,dx
\leq C\|f\|_{p}^{2}\|g\|_{r}+C\Bigl\|\Bigl(\sum|P_{j}f|^{2}\Bigr)^{\frac{1}{2}}\Bigr\|_{p}^{2}\,\|M_{s}g\|_{r}.
\]
既然我们在证明一个乘子定理, 可以假设 $p>2$; 于是存在 $s>1$ 使 $\|M_s(g)\|_r\leq C\|g\|_r$. 总之,
\[
\|T f\|_{p}\leq C\Bigl\|\Bigl(\sum|P_{j}f|^{2}\Bigr)^{\frac{1}{2}}\Bigr\|_{p}\leq C\|f\|_{p}. \qedhere
\]
\end{proof}

\section{评注与推广}
\label{sec:4.3}

\noindent (4.3):\quad 这里陈述的结果与给出的证明取自 H\"ormander~\cite{hormander}; Mikhlin~\cite{mikhlin} 有一个类似的结果. 尽管 $|x|^{-1}$ 在无穷远处不可积, 这些结果在一般维数仍然成立. 这意味着一般的定理需要更多的导数.

\noindent (4.9):\quad 这个结果蕴含 Rademacher 函数是规范正交的, 但它们不构成完备集.

\noindent (4.11):\quad 这里的叙述取自 Stein~\cite{stein-singular}; 原始结果见 Paley~\cite{paley}.

\noindent (4.12):\quad 这里给出的证明仿照了 A.~C\'ordoba 与 C.~Fefferman~\cite{cordoba-feffermanC-weighted} 中的结果, $M_{s}(g)$ 就是在那里引进的. 也有一个使用向量值 Hilbert 变换理论的替代证明; 见 Stein~\cite{stein-singular}.

\noindent (4.13):\quad 这称为 Littlewood--Paley 分解, 得名于圆周上的类似结果. 见 Stein~\cite{stein-square} 的文章.

\noindent (4.14):\quad 这归功于 Marcinkiewicz~\cite{marcinkiewicz-interpolation}; 它在一般维数同样成立. 需要指出的是, 在 $\xi=0$ 处没有任何光滑性条件. 高维的结果在坐标轴上也没有光滑性条件. 这里的证明是 Jodeit 向我们展示的 (未发表).
